\documentclass[11pt]{article}
\usepackage{amsmath,amssymb,amsthm,algorithm,algorithmic}
\usepackage[margin=1in]{geometry}

\newtheorem{theorem}{Theorem}
\newtheorem{lemma}[theorem]{Lemma}
\newtheorem{proposition}[theorem]{Proposition}
\newtheorem{definition}[theorem]{Definition}

\title{Problem 8: Largest Product in a Series}
\author{Project Euler Solutions}
\date{}

\begin{document}
\maketitle

\section{Problem Statement}

Given a 1000-digit number $D = d_0 d_1 \cdots d_{999}$ with $d_i \in \{0, 1, \ldots, 9\}$, find the maximum product of $k = 13$ consecutive digits:
\[
M = \max_{0 \le i \le 987}\; \prod_{j=0}^{12} d_{i+j}.
\]

\section{Mathematical Development}

\subsection{Window Products}

\begin{definition}
For integers $i$ and $k$ with $0 \le i \le N-k$, define the \emph{window product}
\[
P_i = \prod_{j=0}^{k-1} d_{i+j}.
\]
\end{definition}

For this problem, $N = 1000$ and $k = 13$, so there are
\[
N-k+1 = 1000-13+1 = 988
\]
candidate windows.

\begin{theorem}[Exhaustive search is sufficient]\label{thm:exhaustive}
The desired quantity is
\[
M = \max_{0 \le i \le N-k} P_i,
\]
so evaluating all $988$ window products and retaining the largest one yields the correct answer.
\end{theorem}

\begin{proof}
The set
\[
\mathcal{W} = \{P_i : 0 \le i \le N-k\}
\]
is finite. Therefore it has a maximum element. An exhaustive search computes every element of $\mathcal{W}$ and returns the largest value encountered, which is exactly $\max \mathcal{W} = M$.
\end{proof}

\begin{lemma}[Zero windows]
If one of the digits in a length-13 window is zero, then the product of that window is zero.
\end{lemma}

\begin{proof}
A product containing a factor equal to zero is zero.
\end{proof}

\noindent\emph{Remark.} This lemma explains why many windows are automatically non-optimal, but the reference implementation does not require any special optimization: it simply computes all $988$ products directly.

\subsection{Numerical Evaluation}

\begin{proposition}
The maximum product is attained by the 13-digit block
\[
5576689664895,
\]
which begins at index $197$ (using zero-based indexing), and its product is
\[
23\,514\,624\,000.
\]
\end{proposition}

\begin{proof}
Exhaustive evaluation of the $988$ window products shows that the largest one occurs for the block
\[
d_{197} d_{198} \cdots d_{209} = 5576689664895.
\]
Its product is computed directly as
\[
\begin{aligned}
5 \cdot 5 &= 25,\\
25 \cdot 7 &= 175,\\
175 \cdot 6 &= 1050,\\
1050 \cdot 6 &= 6300,\\
6300 \cdot 8 &= 50400,\\
50400 \cdot 9 &= 453600,\\
453600 \cdot 6 &= 2721600,\\
2721600 \cdot 6 &= 16329600,\\
16329600 \cdot 4 &= 65318400,\\
65318400 \cdot 8 &= 522547200,\\
522547200 \cdot 9 &= 4702924800,\\
4702924800 \cdot 5 &= 23514624000.
\end{aligned}
\]
Hence the maximal product is $23\,514\,624\,000$.
\end{proof}

\section{Algorithm}

\begin{algorithm}[H]
\caption{Largest product of $k$ consecutive digits}
\begin{algorithmic}[1]
\STATE $\text{best} \gets 0$
\FOR{$i = 0$ \textbf{to} $N-k$}
    \STATE $P \gets 1$
    \FOR{$j = 0$ \textbf{to} $k-1$}
        \STATE $P \gets P \cdot d_{i+j}$
    \ENDFOR
    \IF{$P > \text{best}$}
        \STATE $\text{best} \gets P$
    \ENDIF
\ENDFOR
\RETURN best
\end{algorithmic}
\end{algorithm}

\begin{theorem}[Algorithm correctness]
The algorithm returns the largest product of $k$ consecutive digits.
\end{theorem}

\begin{proof}
For each index $i$ with $0 \le i \le N-k$, the inner loop computes exactly
\[
P_i = \prod_{j=0}^{k-1} d_{i+j}.
\]
The outer loop visits every admissible starting position exactly once and stores in \texttt{best} the maximum value seen so far. By Theorem~\ref{thm:exhaustive}, the maximum over these values is precisely the desired quantity. Therefore the returned value is correct.
\end{proof}

\section{Complexity Analysis}

\begin{theorem}
The algorithm runs in $\Theta((N-k+1)k)$ time and uses $O(N)$ space in the given implementation.
\end{theorem}

\begin{proof}
There are exactly $N-k+1$ windows, and for each window the inner loop performs exactly $k$ multiplications. Hence the running time is
\[
\Theta((N-k+1)k).
\]
For this problem,
\[
(N-k+1)k = 988 \cdot 13 = 12\,844,
\]
so the computation is tiny. The Python implementation first stores the 1000 digits in a list, which uses $O(N)$ space; aside from that list, only a constant number of scalar variables are maintained.
\end{proof}

\section{Answer}

\[
\boxed{23514624000}
\]

\end{document}
